\section{总结}

\subsection{作品成果}
本作品在一块安路 EG4S20 FPGA 开发板上，实现了一台功能完整、即插即用的数字艺术多媒体画框：
\begin{itemize}
\item \textbf{功能全面}：赛题的全部基础要求与五类拓展功能均已实现，并额外打造了四季动态夜城这一特色展示内容；
\item \textbf{一镜到底}：所有功能集成在同一个固化镜像中，插上 TF 卡、接上显示器、通电即可运行，现场无需电脑和任何调试操作；
\item \textbf{音画一体}：图片、音乐、字幕、转场与动态场景同时工作，切换内容时音乐不断、画面不闪；
\item \textbf{操作简单}：四个按键和三个拨码开关配合中文菜单，即可完成切图、选曲、调节画面和切换场景，普通观众无需说明书也能上手。
\end{itemize}

\subsection{应用效果}
作品把静态图片、背景音乐和随时间变化的数字艺术画面融为一体，观众既能看到清晰的图片内容和中文说明，又能感受到画面随音乐跳动、随季节变换的动态氛围。这种“图片＋音乐＋动态艺术”的组合，使原本单调的信息屏变成具有观赏性的展示终端，更容易吸引观众驻足。

作品不依赖网络、不需要上位机，只要准备好 TF 卡中的图片和音乐即可长期离线运行；更换展示内容时只需替换卡中素材，现场人员无需任何专业知识。基于 FPGA 的纯硬件实现上电即用、运行稳定，适合长时间无人值守的展示环境。

\subsection{应用前景}
\begin{itemize}
\item \textbf{展厅与博物馆}：展品图片配合中文字幕讲解与背景音乐，空闲时以动态夜城营造艺术氛围，提升参观体验；
\item \textbf{商业空间与橱窗}：自动轮播产品图片，动态场景作为吸引视线的数字艺术背景，营造独特的品牌氛围；
\item \textbf{校园与公共信息屏}：轮播通知与安全提示，更新内容只需替换 TF 卡，部署与维护成本低；
\item \textbf{家居数字画框}：播放家庭照片与喜爱的音乐，四季场景提供随时间变化的装饰画面，让墙面“活”起来；
\item \textbf{教学演示}：作为 FPGA 多媒体处理的完整实例，直观展示数字电路如何同时处理图像、声音与交互。
\end{itemize}

\subsection{可扩展之处}
作品采用模块化设计，各部分接口清晰，便于进一步扩展：可以增加新的场景主题，例如节日或城市地标风格；可以更换字幕文案以适配不同展示内容；还可以接入更大的存储或更高的显示分辨率，以容纳更多图片和更精细的画面。新功能可沿用现有的场景控制、图层合成与交互框架，扩展成本较低。

\subsection{作品总结}
本作品以“赛博美学数字艺术多媒体交互画框”为主题，完整实现了赛题要求的 HDMI 多媒体播放系统，并以四季动态夜城展现了 FPGA 在实时数字艺术方面的潜力。全部功能在一块开发板、一个固化镜像中协同运行，开机即用、操作直观。作品表明，仅凭一颗国产 FPGA 芯片，就能构建一台兼具实用性与观赏性的离线音画展示终端，为展陈、商业、校园和家居等场景提供一套完整、可落地的解决方案。
